% !TEX TS-program = xelatex
% !TEX encoding = UTF-8 Unicode
% !Mode:: "TeX:UTF-8"

\documentclass{resume}
\usepackage{zh_CN-Adobefonts_external} % Simplified Chinese Support using external fonts (./fonts/zh_CN-Adobe/)
% \usepackage{NotoSansSC_external}
% \usepackage{NotoSerifCJKsc_external}
% \usepackage{zh_CN-Adobefonts_internal} % Simplified Chinese Support using system fonts
\usepackage{linespacing_fix} % disable extra space before next section
\usepackage{cite}

\begin{document}
\pagenumbering{gobble} % suppress displaying page number

\name{李浩文}

\basicInfo{
  \email{haowenli@smail.nju.edu.cn} \textperiodcentered\ 
  \phone{(+86) 182-6062-0175} 
  % \textperiodcentered\ 
  % \linkedin[billryan8]{https://www.linkedin.com/in/billryan8}
  }
 
\section{\  教育背景}
% \section{\faGraduationCap\  教育背景}
\datedsubsection{\textbf{南京大学}, 南京, 江苏}{2022 -- 至今}
\textit{直博在读}\ 声学, 预计 2027 年 6 月毕业
\datedsubsection{\textbf{南京大学}, 南京, 江苏}{2018 -- 2022}
\textit{本科}\ 声学

\section{\ 项目经历}
% \section{\faUsers\ 实习/项目经历}
\datedsubsection{\textbf{居室降噪项目}}{2022年10月 -- 2024年12月}
\role{华为合作项目}{主要负责人}
% \begin{onehalfspacing}
室内复杂声场环境下，非接触式局部有源降噪系统设计
\begin{itemize}
  \item \textbf{虚拟传感算法：}调研多种常用虚拟传感算法，包括 VM, RM, AF 算法，使用 MATLAB 进行算法实现与仿真，对比分析稳态性能、收敛速度等关键指标
  \item \textbf{实时自适应 ANC 开发：}针对 TI TMS320C6678 开发板，进行 C 语言嵌入式 DSP 开发，实现结合 RM 虚拟传感的多通道实时自适应有源噪声控制
  \item \textbf{改善不同噪声收敛性能：}通过实现时域预白化 FxLMS 算法，以及改进频域归一化 FBFxLMS 算法两种技术路径，解决系统对于有色噪声收敛性能差的问题，改善不同噪声及多噪声源场景下的系统适应性
  \item \textbf{结合人头追踪扩大静区范围：}使用 Intel Realsense 结构光相机，通过 RTMpose 神经网络结合后处理算法，对人头在三维空间的平动坐标和转动角度进行实时定位。ANC 控制器根据定位坐标，实时更新预存滤波器数据，实现静区跟随人头移动
  \item \textbf{空间插值进一步优化降噪性能：}结合人头周围多个建模点信息，利用重心插值算法，拟合当前位置滤波器系数，相比传统的就近切换策略显著提升非建模点位置的降噪量，减小人头移动时降噪量波动，同时避免滤波器切换时的声压跳变，改善体验
  \item \textbf{抗非目标干扰：}系统存在非目标干扰时，通过监测参考信号和误差信号异常，保护系统稳定性
\end{itemize}
% \end{onehalfspacing}

\datedsubsection{\textbf{主动超材料项目}}{2024年5月 -- 2024年12月}
\role{合作课题}{参与项目}
通过主动控制手段，调控材料表面反射波，实现可变的反射系数曲线
\begin{itemize}
  \item 反馈中和算法的 DSP 实现
  \item 子带算法 DSP 实现，从而调控不同子带的反射系数
\end{itemize}

\datedsubsection{\textbf{虚拟声屏障项目}}{2024年5月}
\role{中驰集团合作项目}{参与项目}
在房间窗户布设的虚拟声屏障，阻断噪声的同时保持通风
\begin{itemize}
  \item 物理系统的调试测试
  \item GUI 图形操作界面调试
\end{itemize}

\datedsubsection{\textbf{Hybrid ANC 算法设计}}{2022年6月 -- 2022年10月}
\role{个人项目}{探索研究}
将自适应滤波、卡尔曼滤波等传统算法与神经网络相结合的有源噪声控制算法设计
\begin{itemize}
  \item 调研现有几种 hybrid ANC 策略，包括完全用神经网络全替代自适应滤波、用神经网络辅助选取自适应算法更新步长、用神经网络辅助选取自适应算法更新初值
  \item 复现测试 DNN-FDAF, KalmanNet 等算法
\end{itemize}

\datedsubsection{\textbf{车载有源噪声控制项目}}{2022年1月 -- 2022年6月}
\role{个人项目}{探索研究}
基于实测汽车路噪数据，仿真对比不同 FxLMS 改进算法的性能
\begin{itemize}
  \item 涉及算法包括泄露算法、BCD 算法、NLMS算法、TLMS算法、球面流形优化算法
\end{itemize}

\datedsubsection{\textbf{耳机有源噪声控制项目}}{2021年7月 -- 2021年12月}
\role{个人项目}{探索研究}
基于实测的耳机传递函数，设计有源降噪反馈滤波器
\begin{itemize}
  \item 实验测量不同佩戴情况下耳机次级路径传递函数
  \item 基于内模结构，使用 $H_2/H_\infty$ 方法设计反馈控制滤波器
  \item 使用顺序二次规划、遗传算法等优化算法拟合对应 IIR 滤波器
\end{itemize}

% Reference Test
% \datedsubsection{\textbf{Paper Title\cite{zaharia2012resilient}}}{May. 2015}
% An xxx optimized for xxx\cite{verma2015large}
% \begin{itemize}
%  \item main contribution
% \end{itemize}

\section{\ IT 技能}
% \section{\faCogs\ IT 技能}
% increase linespacing [parsep=0.5ex]
\begin{itemize}[parsep=0.5ex]
  \item 编程语言: C == Python > C++ > Java
  \item 平台: Linux
  \item 开发: xxx
\end{itemize}

\section{\ 获奖情况}
% \section{\faHeartO\ 获奖情况}
\datedline{\textit{第一名}, xxx 比赛}{2013 年6 月}
\datedline{其他奖项}{2015}

\section{\ 其他}
% \section{\faInfo\ 其他}
% increase linespacing [parsep=0.5ex]
\begin{itemize}[parsep=0.5ex]
  \item 技术博客: http://blog.yours.me
  \item GitHub: https://github.com/username
  \item 语言: 英语 - 熟练(TOEFL xxx)
\end{itemize}

% Reference
% \newpage
\renewcommand{\refname}{个人成果}
% % \cite {"Citation Key"}
\nocite{*}    %显示所有文献
% \bibliographystyle{IEEETran}
% \bibliography{mycite}

\begin{thebibliography}{9}
\bibitem{liu_active_2025} Y. Liu, \textbf{H. Li}, H. Zou, Z. Lin, and J. Lu, “Active headrest combined with a depth camera-based ear-positioning system,” The Journal of the Acoustical Society of America, vol. 157, no. 1, pp. 519–526, Jan. 2025, doi: 10.1121/10.0034860.
\bibitem{hu_online_2025} M. Hu, \textbf{H. Li}, J. Lu, H. Zou, and Q. Ma, “An online modeling virtual sensing technique based on kriging interpolation for active noise control,” Mechanical Systems and Signal Processing, vol. 224, p. 112186, Feb. 2025, doi: 10.1016/j.ymssp.2024.112186.
\bibitem{3} 2024 全国声学大会口头汇报“人头移动鲁棒的有源降噪头靠系统设计”
\end{thebibliography}

\end{document}
